% ===========================================================================
%  utente/02-overview/00-capitolo.tex
% ===========================================================================
\chapter{Overview}\label{cap:overview}

È la pagina che si apre per prima. Risponde a una domanda sola: come va
tutto, adesso.

\figuraSchermata{1.0}{figure/schermate/overview.png}
  {La schermata iniziale. Quattro riquadri: il grafico dell'ensemble, le
   metriche, l'elenco degli stream attivi, e la scorciatoia verso lo Studio.}
  {overview}

\section{Il grafico in alto a sinistra}\label{sez:overview-grafico}

Mostra l'andamento del \voce{Project} scelto: la combinazione di più
strategie, al netto dei costi stimati, sui dati che la strategia
\emph{non} aveva visto mentre veniva costruita.

Finché non si è salvato nessun project, al suo posto compare un invito a
crearne uno. Non è un errore: è che non c'è ancora niente da disegnare.

\section{Metrics OOS}\label{sez:overview-metriche}

Sei numeri, e vanno letti insieme. \voce{OOS} sta per \emph{out of sample}:
calcolati sui dati che la strategia non conosceva.

\begin{center}
\begin{tabular}{@{}L{32mm}L{106mm}@{}}
\toprule
\textbf{Metrica} & \textbf{Che cosa dice} \\
\midrule
\rigaTenue
\voce{Sharpe ratio} & Quanto rende rispetto a quanto oscilla. È il numero
  più usato per confrontare due strategie, e il più facile da far sembrare
  bello guardando i dati sbagliati. \\
\voce{CAGR net} & Il rendimento annuo composto, già al netto dei costi
  stimati. \\
\rigaTenue
\voce{Max drawdown} & La perdita più grande dal massimo precedente. È la
  domanda «quanto avrei dovuto sopportare prima di recuperare». \\
\voce{Ann. volatility} & Quanto oscilla in un anno. \\
\rigaTenue
\voce{Alpha vs SPY ann.} & Quanto rende \emph{in più} del mercato, tolta la
  parte che si spiega con il mercato stesso. Una strategia con un buon
  rendimento e alfa vicino a zero sta comprando il mercato con altri nomi. \\
\voce{Win rate} & La percentuale di giorni in guadagno. Da solo dice poco:
  si può vincere nove volte su dieci e perdere comunque. \\
\bottomrule
\end{tabular}
\end{center}

\begin{perche}{Quali numeri guardare per primi}
\voce{Max drawdown} e \voce{Alpha vs SPY}. Il primo perché è la cifra che
decide se una strategia è sopportabile: uno Sharpe eccellente con un
drawdown del \num{45}\,\% è una strategia che nessuno tiene fino in fondo.
Il secondo perché distingue una strategia che fa qualcosa da una che sta
seguendo il mercato con passi più lunghi.
\end{perche}

\begin{attenzione}{Un backtest non è una previsione}
Dice come una strategia si \emph{sarebbe} comportata sui dati del passato,
con i costi stimati. Non dice come si comporterà. Il capitolo
\vediCap{streams} spiega quali segnali distinguono una strategia che ha
funzionato da una che ha avuto fortuna.
\end{attenzione}

\section{Active streams}\label{sez:overview-elenco}

L'elenco di tutti i risultati disponibili, ordinato. Le colonne:

\begin{description}[leftmargin = 1.4em, style = nextline, itemsep = 4pt]
  \item[\voce{Sharpe}] Come sopra.
  \item[\voce{OOS Return}] Il rendimento fuori campione. Verde sopra zero,
    rosso sotto.
  \item[\voce{Corr}] La correlazione con gli altri stream. È la colonna che
    conta di più quando si combinano: due strategie molto correlate, messe
    insieme, non diversificano niente --- fanno la stessa scommessa due volte.
  \item[\voce{Watch}] La stellina mette lo stream fra i preferiti.
\end{description}

\noindent
I tre menu in alto all'elenco scelgono il periodo (\voce{OOS} o tutto),
l'ordinamento e il project di riferimento.

\section{Tutorials}\label{sez:overview-tutorial}

In basso a destra, quattro passi che si spuntano man mano che si fanno:
costruire il primo stream, validarlo, combinarlo, metterlo in produzione. È
il percorso che questo manuale segue nei capitoli successivi.
